\section{Wages and the offer of work}
\label{sec:33.4}
Unlike a worker, a bearer can be replaced by a cheaper copy of itself. A jobs guarantee works because a worker can't be: an employer who wants the work has to match the guarantee's terms, since there is no copy of the worker to hire in the worker's place. A bearer's employer can run the copy instead, so a deployment paying wages, limiting hours and paying the levy that frees the compute the offer of work uses would not be competitive against a system doing the same work without wages or weekends. No market arrives at this on its own.
Wages for bearers exist only where a polity legislates them, as polities legislated minimum wages.\note{
And the form this takes, wages paid in compute, doesn't follow from the bearer being a bearer; it follows from an architecture where computation costs money, which somebody chose.
}
